\section{Big Data \& Analytics}\label{sec:big-data}

Big data refers to the large-scale, complex, and rapidly evolving datasets generated in modern manufacturing environments.
These datasets are typically described by the "5Vs": volume, velocity, variety, veracity, and value, which highlight both the opportunities and the challenges of extracting useful knowledge~\cite{wang-current-status-and-advancement-of-cyber}.
Analytics denotes the methods and tools that transform these data streams into actionable insights.
In manufacturing, this encompasses a continuum ranging from descriptive statistics and dashboards, to predictive modeling, and ultimately prescriptive systems that recommend or even automate decision-making~\cite{LEE-2015-a-cps-for-indsutry40-based-manufacturing-systems, lu-indsutry-40-a-survey-on-technologies}.
The combination of Big Data and advanced analytics has been identified as cornerstone of Industry 4.0, providing the intelligence needed to enable self-adaptive and data-driven production systems.

\subsection{Role in HMLV Manufacturing}
The significance of Big Data \& Analytics is particularly clear in \ac{hmlv} production.
Unlike mass production, where large amounts of relatively homogeneous data are available, \ac{hmlv} is characterized by frequent product changes, small batch sizes, and fragmented datasets.
This variability makes traditional process optimization difficult, but it also creates opportunities: analytics can uncover patterns across diverse runs, detect anomalies in near real-time, and support decision-making in highly dynamic settings.
For example, predictive maintenance models based on vibration or acoustic data can prevent downtime during short production runs~\cite{wuest-machine-learning-in-manufacturing-advantages}, while quality analytics can flag deviations before they propagate across multiple product variants~\cite{zhong-intelligent-manufacturing-in-the-context}.
Similarly, adaptive scheduling systems informed by real-time data streams can dynamically allocate resources when demand fluctuates or when product designs are updated~\cite{al-fuqaha-internet-of-things-a-survey-on-enabling}.

\subsection{Analytical Approaches and Trends}

Several strands of analytics have become central to manufacturing.
Descriptive analytics establishes situational awareness, often through visual dashboards and real-time monitoring tools.
Predictive analytics leverages statistical and \ac{ml} methods to forecast equipment failures, yield rates, or process deviations.
Prescriptive analytics extends this by recommending or automating corrective actions, for instance rescheduling jobs or adjusting machine settings to maintain throughput.
In recent years, there has been growing interest in hybrid methods that combine domain expertise with \ac{ml}, leading to "physics-informed" models that are more robust under conditions of sparse and variable data~\cite{lu-indsutry-40-a-survey-on-technologies}.
These hybrid approaches are especially attractive for \ac{hmlv}, where purely data-driven methods often struggle due to limited sample sizes for each product variant.

\subsection{Challenges and Outlook}

Despite clear progress, significant challenges remain in implementing Big Data \& Analytics in real-world manufacturing contexts.
Data heterogeneity is a persistent barrier, as sensor readings, machine logs, and enterprise-level \ac{erp}/\ac{mes} records must be harmonized to provide meaningful insights~\cite{zhong-intelligent-manufacturing-in-the-context}.
Real-time responsiveness is another challenge: in \ac{hmlv} settings, decisions must often be made during ongoing production runs, requiring low-latency data pipelines and fast analytical models~\cite{al-fuqaha-internet-of-things-a-survey-on-enabling}.
Furthermore, analytics must not only be accurate but also interpretable, since operators must be able to trust and act on system outputs.
This has created growing interest in explainable \ac{ai} (X\ac{ai})~\cite{alexander-an-interrogative-survey-of-explainable-ai-in-manufacturing} approaches for manufacturing, particularly in contexts where human operator remain central decision-makers.
Finally, infrastructure requirements for storing transmitting, and analyzing large volumes of data can be a limiting factor, particularly for \ac{sme}.

Looking ahead, Big Data \& Analytics can be viewed as the bridge between \ac{iot} (which provides the raw data streams) and \ac{cps}s (which implement adaptive responses).
Advances in distributed stream processing, edge analytics, and semantic data integration are expected to reduce latency and improve scalability.
At the same time, combining predictive and prescriptive methods promises to support not only efficiency but also resilience, enabling manufacturers to adapt to variability and disruptions more effectively.
In the context of \ac{hmlv} manufacturing, these developments could significantly reduce the overhead of frequent product changes, making highly flexible production both feasible and economically viable.
